% =====================================================================
\section{Les questions probables, et quoi répondre}

\subsection{Architecture}

\textbf{\enquote{Pourquoi des servlets et des JSP, et pas Spring Boot ?}}\\
Le projet reste sur la pile \textbf{Jakarta EE standard} : les servlets, les
EJB et JPA font partie de la plateforme, et WildFly les fournit. Il n'y a donc
aucune dépendance de \emph{framework} à gérer dans le \texttt{pom.xml}, et
l'application se déploie sur n'importe quel serveur compatible. Spring Boot
aurait apporté de la configuration automatique, au prix d'une pile
supplémentaire à maîtriser et à justifier.

\textbf{\enquote{Qu'est-ce qui garantit la séparation des couches ?}}\\
Trois choses vérifiables dans le code : aucun servlet n'injecte
d'\texttt{EntityManager} --- ils n'ont que des \texttt{@EJB} ; aucun service ne
manipule \texttt{HttpServletRequest} ; aucune JSP ne contient de scriptlet
Java. Chaque couche ne connaît que la suivante.

\textbf{\enquote{Pourquoi les JSP sont-elles sous \texttt{WEB-INF} ?}}\\
Parce que le conteneur en interdit l'accès par URL directe. On ne peut y entrer
que par un \texttt{forward} depuis un servlet, donc après le contrôle de
session et de rôle. Les deux exceptions, \texttt{index.jsp} et
\texttt{forgot-password.jsp}, sont volontaires : elles doivent être
atteignables sans être connecté.

\subsection{Sécurité}

\textbf{\enquote{Où est contrôlé le rôle ?}}\\
Au début de \texttt{doGet} et \texttt{doPost} de chaque servlet concerné. Le
menu masque aussi les entrées interdites, mais ce n'est \textbf{qu'un
confort} : la vue est contournable, le contrôle serveur ne l'est pas.

\textbf{\enquote{Et si quelqu'un forge une requête HTTP ?}}\\
Sur le circuit de validation, trois barrières se superposent : la vue
n'affiche le bouton que pour le bon rôle et le bon état ; le servlet croise
l'action et le rôle et rejette toute combinaison non prévue ; le service
vérifie l'état de départ de la demande avant d'agir. Les deux dernières sont
côté serveur.

\textbf{\enquote{Comment vous protégez-vous de l'injection SQL ?}}\\
Aucune requête n'est construite par concaténation. Tout passe par JPQL avec
\texttt{setParameter}, donc par des variables liées : la valeur n'est jamais
interprétée comme du code SQL.

\textbf{\enquote{Le mot de passe est-il protégé ?}} --- \emph{à annoncer
soi-même plutôt que de le subir}\\
Non : il est comparé en clair dans \texttt{SecurityService.authenticate}. C'est
la limite principale du projet. La correction est localisée : stocker un
condensat salé (BCrypt, Argon2), charger l'utilisateur par son seul
identifiant, puis comparer en Java. Une seule méthode et une migration de la
colonne.

\subsection{JPA et transactions}

\textbf{\enquote{Qui ouvre et ferme les transactions ?}}\\
Le conteneur. L'unité de persistance est en \texttt{transaction-type="JTA"} et
les services sont des EJB : une transaction est ouverte à l'entrée de chaque
méthode et validée à la sortie, ou annulée si une exception s'échappe. C'est
la raison de la présence de \texttt{CMTTxInterceptor} dans les traces.

\textbf{\enquote{Pourquoi \texttt{hbm2ddl.auto = none} ?}}\\
Pour qu'Hibernate ne touche jamais au schéma. Les tables sont créées par les
scripts de \texttt{sql/}, et la base fait autorité. Un \texttt{update}
automatique en production pourrait modifier une structure existante sans
contrôle.

\textbf{\enquote{Différence entre \texttt{persist} et \texttt{merge} ?}}\\
\texttt{persist} enregistre une entité nouvelle ; \texttt{merge} réintègre une
entité détachée et renvoie l'instance gérée. Le projet applique le motif
classique : \texttt{if (entite.getId() == null) persist else merge}.

\subsection{Front}

\textbf{\enquote{Pourquoi pas de framework front ?}}\\
Les pages sont des listes et des formulaires, rendus par le serveur : un
client riche n'apporterait rien et ajouterait une étape de construction, un
gestionnaire de paquets et des dépendances à suivre. Deux bibliothèques
suffisent : ApexCharts pour les graphiques, Lucide pour les icônes.

\textbf{\enquote{Comment les données arrivent-elles dans les graphiques ?}}\\
Le servlet pose des chaînes déjà formatées avec \texttt{setAttribute}, la JSP
les insère dans la configuration ApexCharts. Pas de sérialisation JSON ni
d'appel asynchrone. La limite est qu'un libellé contenant une apostrophe
casserait le script : pour des données saisies par l'utilisateur, il faudrait
du JSON échappé.

\subsection{Décisionnel}

\textbf{\enquote{Pourquoi des vues \texttt{V\_BI\_*} plutôt que les tables ?}}\\
Pour poser une \textbf{couche sémantique} : les colonnes techniques deviennent
des indicateurs métier, les clés deviennent des libellés, les valeurs nulles
sont ramenées à zéro et les jointures sont faites une fois pour toutes. Power
BI consomme du sens, pas de la structure. Et si le schéma évolue, seule la vue
change, pas le rapport.

\textbf{\enquote{Pourquoi une chaîne ETL en plus de Power BI ?}}\\
Parce que les deux ne servent pas au même usage. Power BI restitue au niveau du
détail ; la segmentation a besoin d'un \textbf{datamart agrégé par agence},
avec un contrôle de qualité et une sortie stable, reproductible d'une exécution
à l'autre. L'ETL produit ce jeu, et le clustering le consomme.

\textbf{\enquote{Mesure ou colonne calculée ?}}\\
Une colonne calculée est évaluée à l'actualisation et occupe de la mémoire :
elle sert à \emph{filtrer ou grouper}. Une mesure est évaluée à l'affichage,
dans le contexte des filtres actifs : elle sert à \emph{agréger}. Dans le
rapport, \texttt{Rôle} est une colonne, \texttt{Taux Présence} une mesure.

\textbf{\enquote{Pourquoi K-Means et pas un autre modèle ?}}\\
Parce que c'est celui qui obtient le meilleur score de silhouette (0,604)
\textbf{sous un protocole équitable} : les quatre modèles sont notés sur la
totalité des agences, et les points que DBSCAN refuse de classer sont comptés
comme un groupe. Sans cette précaution, un modèle serait avantagé pour avoir
écarté les cas difficiles.

\textbf{\enquote{Pourquoi écarter le siège de la segmentation ?}}\\
540 employés contre une dizaine pour une agence : c'est une structure
centrale, pas une agence. Le laisser reviendrait à créer un groupe d'un seul
point et à écraser les distances entre toutes les autres. Il reste dans le
datamart, il est seulement retiré de l'analyse.

% =====================================================================
\section{Points ouverts et limites}

Les annoncer soi-même vaut mieux que de les subir. Ils ont été relevés par un
inventaire du code fait pour ce dossier.

\noindent\begin{tabularx}{\textwidth}{|>{\columncolor{btkpale}\bfseries}L{3.6cm}|X|}
\hline
Mot de passe en clair & Comparé tel quel dans \texttt{SecurityService.authenticate}. Correction : condensat salé, comparaison en Java.\\
\hline
Contrôle de rôle partiel & Quatre servlets ne vérifient que l'authentification, sans le rôle : \texttt{BUtilisateurServlet} (\texttt{/employes}), \texttt{BObjectifServlet} (\texttt{/objectifs-btk}), \texttt{ObjectifServlet} (\texttt{/objectifs}) et \texttt{ClientServlet}. Ce sont des écrans de consultation, mais l'application est donc plus permissive que la matrice des droits.\\
\hline
Page gestionnaires & \texttt{GestionnaireServlet} transfère vers \texttt{/WEB-INF/gestionnaires.jsp} alors que le fichier s'appelle \texttt{gestionnaire.jsp} : la page échoue. Elle n'est pas dans le menu. Correction : un mot dans le chemin du \texttt{forward}.\\
\hline
Reliquat historique & \texttt{ClientServlet} (\texttt{/clients-old}) et \texttt{client.jsp} ne sont atteignables par aucun lien du menu ; \texttt{client.jsp} pointe vers \texttt{/clients}, qui n'est associé à aucun servlet.\\
\hline
Pas de tests automatisés & Le projet n'a pas de tests unitaires. Les services étant des EJB sans état, ils s'y prêteraient bien (JUnit + base en mémoire).\\
\hline
Graphiques et échappement & Les séries ApexCharts sont injectées comme chaînes ; un libellé contenant une apostrophe casserait le script.\\
\hline
\end{tabularx}

\vspace{0.8em}
\begin{retenir}
\textbf{Comment présenter une limite.} La formuler, en donner la cause, puis
la correction et son coût. \enquote{Le mot de passe est comparé en clair ;
c'est un choix de périmètre, pas un oubli ; la correction tient dans une
méthode et une migration de colonne.} Une limite reconnue et chiffrée est un
signe de maîtrise ; une limite découverte par le jury ne l'est pas.
\end{retenir}

% =====================================================================
\section{Récapitulatif : la carte des URL}

\noindent\begin{longtable}{|>{\ttfamily\small}L{3.3cm}|>{\small}L{3.9cm}|>{\small}L{3.1cm}|>{\small}L{4.1cm}|}
\hline
\rowcolor{btkpale}
\textmd{\textbf{URL}} & \textbf{Servlet} & \textbf{Rôles} & \textbf{Vue}\\
\hline
\endhead
/login & LoginServlet & ouverte & index.jsp\\
\hline
/logout & LogoutServlet & connecté & ---\\
\hline
/forgot-password & ForgotPasswordServlet & ouverte & forgot-password.jsp\\
\hline
/home & HomeServlet & connecté & home.jsp\\
\hline
/employes & BUtilisateurServlet & connecté & employes.jsp\\
\hline
/clients-btk & BClientServlet & ADMIN & clients-btk.jsp\\
\hline
/banques & BanqueServlet & ADMIN & banques.jsp\\
\hline
/agences & AgenceServlet & ADMIN & agences.jsp\\
\hline
/objectifs-btk & BObjectifServlet & connecté & objectifs-btk.jsp\\
\hline
/objectifs & ObjectifServlet & connecté & objectifs.jsp\\
\hline
/pointage & PointageServlet & ADMIN (saisie manuelle) & pointage.jsp\\
\hline
/ajouter-employe & AjouterEmployeServlet & ADMIN & add-effectif.jsp\\
\hline
/demandes-agence & DemandeAgenceServlet & ADMIN, DIRECTEUR & demandes-agence.jsp\\
\hline
/demandes-admin & DemandeAgenceAdminServlet & ADMIN & demandes-admin.jsp\\
\hline
/gestion-demandes & GestionModificationsServlet & ADMIN, DIRECTEUR & gestion-demandes.jsp\\
\hline
/mes-demandes & MesDemandesServlet & connecté & profil.jsp\\
\hline
/gestionnaires & GestionnaireServlet & ADMIN & \textit{transfert erroné}\\
\hline
/clients-old & ClientServlet & connecté & \textit{reliquat}\\
\hline
\end{longtable}
